\documentclass[a4paper, 11pt, oneside]{article}
\usepackage[a4paper,margin=1.5cm]{geometry}
\usepackage[utf8]{inputenc}
\usepackage[french]{babel}
\usepackage[T1]{fontenc}
\usepackage[fleqn]{amsmath}
\usepackage{listings}
\usepackage{color}
\usepackage[dvipsnames]{xcolor}
\usepackage{fullpage}
\usepackage{tikz}
\usepackage{bm}

\newcommand\bigforall{\mbox{\Large $\mathsurround0pt\forall$}}

\title{Logic for Computer Science Project : Report}
\author{Kuckart Marco}
\date{}

\begin{document}
\maketitle

\section{Question 1}

For this question let $p_{i,j}$ be the proposition ``the cell $i$ is reached with
$j$ moves''.

First, all cells must be visited: $\forall i \bigvee_j p_{i,j}$

It is obvious that given the definition of $p_{i,j}$, we have the following
constraint: A cell cannot be reached in two different numbers of moves.

So we have
$$\forall i \left(\forall j,k,j\neq k , \lnot(p_{i,j} \land p_{i,k}) \leftrightarrow (\lnot p_{i,j} \lor \lnot p_{i,k})\right)$$

Another intuitive constraint is that two different cells cannot be reached in the
same number of moves.

So we also have
$$\forall i \left(\forall j,k,j\neq k , \lnot(p_{j,i} \land p_{k,i}) \leftrightarrow (\lnot p_{j,i} \lor \lnot p_{k,i})\right)$$

All possible numbers of moves have to be assigned, but this is a logical consequence
of the previous statements.

Let now $R_i$ denote the set of all cells reachable from the cell $i$ ($\# R_i \leq 8$).

Given that they are the only possible moves, we have $\forall i, j$
$$p_{i,j} \Rightarrow \bigvee_{r\in R_i} p_{r,j+1}
\leftrightarrow
\lnot p_{i,j} \bigvee \left(\bigvee_{r\in R_i} p_{r,j+1}\right)$$

These clauses are sufficient to solve the problem.

\section{Question 2}

My first attempt to encode the problem in a more efficient way was the following:
let $s_{i,j}$ be the proposition ``cell $i$ is followed by cell $j$ in the path''.
This encoding has the advantage that the value of most variables can be set to false:
for each cell $i$, there are only at most 8 cells that are valid for $s_{i,j}$, the
others can be set to false.
Then, the solver will have many clauses that are just literals and will therefore
be able to perform a large number of propagations very quickly.

However, I could not find a way to encode the fact that the knight must visit each cell.
I therefore decided to use this new encoding together with the previous variables $p_{i,j}$,
and to add clauses linking the two notions together.
I also found that a few variables can have fixed values, and I removed
all the redundant clauses: $(p \lor q) \land (q \lor p) \leftrightarrow (p \lor q)$.

With all these improvements, the solver is able to solve the problem
for an $8 \times 8$ chessboard about twelve times faster than before.


\section{Question 3}

To force the model to compute a new model, only one additional clause is needed: the
disjunction of the negations of all the position variables that were true in the
previous model.
This constraint means that \textbf{at least one of the cell un the path must be in a
different position than before} and that's it.

\section{Question 4}

Previously, to obtain a new model, the added clause was the disjunction of the
negations of all the position variables $p_{i,j}$ (cf. Question 1) that were true
in the previous model. The idea remains the same, but for each symmetry, we add
the disjunction of the variables $p_{i', j}$ where $i'$ is the cell obtained by
applying the symmetry to cell $i$.

\section{Question 5}

Since the solution is a path, one can visualize the multiple solutions for a path
starting at $(i_0, j_0)$ as a rooted tree.

\begin{figure}[h]
    \centering
    
    \begin{tikzpicture}[scale=0.6]

        \foreach \x in {0,...,11} {
            \foreach \y in {0,...,11} {
                \draw[fill=black, draw=black] (\x, -\y) circle (2pt);
            }
        }

        \foreach \i in {0,...,11} {
            \node at (\i, 0.5) {$\i$}; % x labels
            \node at (-0.5, -\i) {$\i$}; % y labels
        }

        \node[font=\bfseries] at (5.5, 1.2) {
            Cells
        };

        \node[font=\bfseries, rotate=90] at (-1.5, -5.5) {
            Position
        };

        \draw[ultra thick, blue!70] (4, 0) -- (2, -1);
        \draw[ultra thick, blue!70] (2, -1) -- (11, -2);
        \draw[ultra thick, blue!70] (11, -2) -- (5, -3);
        \draw[ultra thick, blue!70] (5, -3) -- (3, -4);
        \draw[ultra thick, blue!70] (3, -4) -- (10, -5);
        \draw[ultra thick, blue!70] (10, -5) -- (1, -6);
        \draw[very thick, MidnightBlue!70] (1, -6) -- (8, -7);
        \draw[very thick, MidnightBlue!70] (8, -7) -- (6, -8);
        \draw[very thick, MidnightBlue!70] (6, -8) -- (0, -9);
        \draw[very thick, MidnightBlue!70] (0, -9) -- (9, -10);
        \draw[very thick, MidnightBlue!70] (9, -10) -- (7, -11);

        \draw[ultra thick, red!70] (4, 0) -- (10, -1);
        \draw[ultra thick, red!70] (10, -1) -- (3, -2);
        \draw[ultra thick, red!70] (3, -2) -- (5, -3);
        \draw[ultra thick, red!70] (5, -3) -- (11, -4);
        \draw[ultra thick, red!70] (11, -4) -- (2, -5);
        \draw[ultra thick, red!70] (2, -5) -- (9, -6);
        \draw[very thick, magenta!70] (9, -6) -- (7, -7);
        \draw[very thick, magenta!70] (7, -7) -- (1, -8);
        \draw[very thick, magenta!70] (1, -8) -- (8, -9);
        \draw[very thick, magenta!70] (8, -9) -- (6, -10);
        \draw[very thick, magenta!70] (6, -10) -- (0, -11);

        \draw[very thick, Turquoise!70] (1, -6) -- (7, -7);
        \draw[very thick, Turquoise!70] (7, -7) -- (9, -8);
        \draw[very thick, Turquoise!70] (9, -8) -- (0, -9);
        \draw[very thick, Turquoise!70] (0, -9) -- (6, -10);
        \draw[very thick, Turquoise!70] (6, -10) -- (8, -11);

        \draw[very thick, orange!70] (9, -6) -- (0, -7);
        \draw[very thick, orange!70] (0, -7) -- (6, -8);
        \draw[very thick, orange!70] (6, -8) -- (8, -9);
        \draw[very thick, orange!70] (8, -9) -- (1, -10);
        \draw[very thick, orange!70] (1, -10) -- (7, -11);

    \end{tikzpicture}


    \caption{Decision tree representing all possible paths for a $3\times 4$ chessboard
    with $i_0 = 1$ and $j_0 = 0$}\label{tree}
\end{figure}

Intuitively, to obtain a list of constraints that ensures a unique solution, my
algorithm enumerates all valid paths using the SAT solver and compares them iteratively.
When several paths differ at a given position, one possible cell is chosen arbitrarily,
and a corresponding constraint is added to keep only the consistent solutions.
This process is repeated until only one path remains.
\\

\noindent\textbf{Note:} One possible output of the algorithm (in the same situation
as in Figure~\ref{tree}) is \texttt{[(1, 2, 2), (7, 1, 3)]} 
As can be seen in Figure~\ref{tree}, this indeed corresponds to a constraint that
ensures a unique solution, the red/pink path.
(The cell $(2,2)$ corresponds to column 10, and the cell $(1,3)$ corresponds to
column 7 in the figure.)


\end{document}